\chapter{Метод получения оптимальной траектории}

В приложении приведены основные формулы и переходы, необходимые для
понимания того, как из задачи оптимального управления получается
оптимальная траектория. Подробные промежуточные выкладки, численные
детали и сравнение с методом продолжения остаются за пределами
настоящего приложения.

\section*{1. Постановка задачи в декартовых координатах}

Движение космического аппарата в задаче двух тел с реактивным
ускорением задаётся системой
\begin{equation}
\label{eq:appA_cartesian_motion}
\dot{\mathbf r}=\mathbf v,\qquad
\dot{\mathbf v}=-\frac{\mu}{r^3}\mathbf r+\mathbf a.
\end{equation}

При идеально регулируемой тяге и ограничении на доступную мощность
оптимизация сводится к минимизации квадратичного функционала
\begin{equation}
\label{eq:appA_J}
J=\frac12\int_{t_0}^{t_f}\Vert\mathbf a(t)\Vert^2\,dt.
\end{equation}

Функция Гамильтона--Понтрягина для этой задачи имеет вид
\begin{equation}
\label{eq:appA_H}
H=-\frac{1}{2}\Vert\mathbf a\Vert^2+
\mathbf p_r\cdot \mathbf v
+\mathbf p_v\cdot\left(-\frac{\mu}{r^3}\mathbf r+\mathbf a\right).
\end{equation}

Из условия стационарности
\(
\partial H/\partial \mathbf a=0
\)
получаем оптимальное управление
\begin{equation}
\label{eq:appA_aopt}
\mathbf a=\mathbf p_v.
\end{equation}

Тем самым исходная задача оптимального управления сводится к
двухточечной краевой задаче для фазовых и сопряжённых переменных:
\begin{equation}
\label{eq:appA_cartesian_tpbf}
\dot{\mathbf r}=\mathbf v,\qquad
\dot{\mathbf v}=-\frac{\mu}{r^3}\mathbf r+\mathbf p_v,\qquad
\dot{\mathbf p}_r=-\frac{\partial H}{\partial \mathbf r},\qquad
\dot{\mathbf p}_v=-\frac{\partial H}{\partial \mathbf v}.
\end{equation}

В задаче встречи на правом конце требуется выполнение условий
\begin{equation}
\label{eq:appA_rendezvous}
\mathbf r(t_f)=\mathbf r_f,\qquad
\mathbf v(t_f)=\mathbf v_f.
\end{equation}

Таким образом, оптимальная траектория полностью определяется начальными
значениями сопряжённых переменных.

\section*{2. Регуляризация: преобразования Сундмана и Кустаанхеймо--Штифеля}

Чтобы уменьшить чувствительность задачи, вводится фиктивное время $s$
по преобразованию Сундмана
\begin{equation}
\label{eq:appA_sundman}
dt=r\sqrt{\frac{a}{\mu}}\,ds,
\end{equation}
а также временная компонента $\tau_{\mathrm{KS}}$, связанная с физическим временем
соотношением
\begin{equation}
\label{eq:appA_tau}
t=\tau_{\mathrm{KS}}-
\frac{1}{\sqrt{-2\mathcal E}}\frac{\mathbf r\cdot \mathbf r'}{r}.
\end{equation}

Далее используется KS-преобразование, переводящее трёхмерный
радиус-вектор в четырёхмерные параметрические переменные:
\begin{equation}
\label{eq:appA_KS}
\begin{pmatrix}
x\\ y\\ z\\ 0
\end{pmatrix}
=
L(\mathbf u)\mathbf u,
\qquad
r=\mathbf u\cdot\mathbf u,
\end{equation}
где
\begin{equation}
\label{eq:appA_Lu}
L(\mathbf u)=
\begin{pmatrix}
u_1 & -u_2 & -u_3 & u_4\\
u_2 &  u_1 & -u_4 & -u_3\\
u_3 &  u_4 &  u_1 &  u_2\\
u_4 & -u_3 &  u_2 & -u_1
\end{pmatrix},
\end{equation}
а допустимые параметрические состояния удовлетворяют билинейному
соотношению
\begin{equation}
\label{eq:appA_bilinear}
\mathbf u_b\cdot \mathbf u=0,
\qquad
\mathbf u_b=(u_4,-u_3,u_2,-u_1).
\end{equation}

После регуляризации уравнения движения принимают вид
\begin{equation}
\label{eq:appA_ks_system}
\mathbf u'=\mathbf w,\qquad
\mathbf w'=-\frac{\mathbf u}{4}
-\frac{\mathbf u\cdot\mathbf u}{4\mathcal E}L^T(\mathbf u)\mathbf a
-\frac{\mathcal E'}{2\mathcal E}\mathbf w,
\end{equation}
\begin{equation}
\label{eq:appA_Eprime}
\mathcal E'=2\mathbf w\cdot L^T(\mathbf u)\mathbf a,
\end{equation}
\begin{equation}
\label{eq:appA_tauprime}
\tau_{\mathrm{KS}}'=\frac{1}{(-2\mathcal E)^{3/2}}
\left(
\mu+4(\mathbf u\cdot\mathbf w)\mathcal E'
+(\mathbf u\cdot\mathbf u)\,\mathbf u\cdot L^T(\mathbf u)\mathbf a
\right).
\end{equation}

В этих переменных задача приобретает форму регуляризованного
возмущённого гармонического осциллятора, а потому становится удобнее
для численного поиска начального приближения.

\section*{3. Оптимальное управление в регулярных переменных}

В регулярных переменных функционал записывается уже по фиктивному
времени:
\begin{equation}
\label{eq:appA_Jks}
J=\frac12\int_{s_0}^{s_f}
\Vert\mathbf a(t)\Vert^2\,
\frac{\Vert\mathbf u(s)\Vert^2}{\sqrt{-2\mathcal E(s)}}\,ds.
\end{equation}

При этом к управлению добавляется смешанное ограничение
\begin{equation}
\label{eq:appA_mixed_constraint}
g(s,\mathbf u,\mathbf a)=\mathbf u_b\cdot L^T(\mathbf u)\mathbf a=0,
\end{equation}
исключающее нефизические траектории.

После применения принципа максимума Понтрягина и введения
вспомогательного вектора $\widetilde{\boldsymbol\Lambda}$ оптимальное управление
получается в компактной форме
\begin{equation}
\label{eq:appA_aopt_ks}
\mathbf a=\widetilde{\boldsymbol\Lambda}\,
\frac{\sqrt{-2\mathcal E}}{\Vert\mathbf u\Vert^2},
\qquad
(\mathbf a)_4=0.
\end{equation}

Именно эта формула замыкает задачу в регулярных переменных: после неё
система для фазовых и сопряжённых переменных снова становится
замкнутой, а искомыми параметрами остаются начальные значения
сопряжённых переменных.

\section*{4. Как из этого получается оптимальная траектория}

Практически метод работает так.

\begin{enumerate}
  \item Задаются начальные и конечные условия перелёта в физических
  переменных.
  \item Выполняется переход к регулярным переменным
  $(\mathbf u,\mathbf w,\mathcal E,\tau_{\mathrm{KS}})$.
  \item По принципу максимума Понтрягина записывается оптимальное
  управление в форме \eqref{eq:appA_aopt_ks}.
  \item Тем самым задача сводится к краевой задаче по начальным
  сопряжённым переменным.
  \item Краевая задача далее редуцируется к задаче оптимизации по
  невязкам на правом конце траектории.
\end{enumerate}

Таким образом, оптимальная траектория получается не прямым подбором
управления вдоль всего перелёта, а через поиск такого начального вектора
сопряжённых переменных, при котором решение соответствующей краевой
задачи удовлетворяет терминальным условиям встречи.

\section*{5. Связь с основной диссертацией}

Материал настоящего приложения нужен в диссертации только как краткая
опора для понимания двух идей:
\begin{enumerate}
  \item оптимальная траектория получается из принципа максимума
  Понтрягина и сводится к краевой задаче;
  \item выбор удобного набора переменных и преобразований может резко
  повлиять на численную устойчивость поиска решения.
\end{enumerate}

Предварительные исследования автора были связаны с получением
оптимальной траектории через краевую задачу и специальные
преобразования переменных. В настоящей диссертации центральным
объектом является уже не единичная траектория, а многообразие таких
решений и его аппроксимация.
